\documentclass[conference]{IEEEtran}
\IEEEoverridecommandlockouts
% The preceding line is only needed to identify funding in the first footnote. If that is unneeded, please comment it out.
\usepackage{cite}
\usepackage{amsmath,amssymb,amsfonts}
\usepackage{algorithmic}
\usepackage{graphicx}
\usepackage{textcomp}
\usepackage{xcolor}
\def\BibTeX{{\rm B\kern-.05em{\sc i\kern-.025em b}\kern-.08em
    T\kern-.1667em\lower.7ex\hbox{E}\kern-.125emX}}
\begin{document}

\title{
\thanks{Identify applicable funding agency here. If none, delete this.}
}

\author{\IEEEauthorblockN{1\textsuperscript{st} Given Name Surname}
\IEEEauthorblockA{\textit{dept. name of organization (of Aff.)} \\
\textit{name of organization (of Aff.)}\\
City, Country \\
email address or ORCID}
\and
\IEEEauthorblockN{2\textsuperscript{nd} Given Name Surname}
\IEEEauthorblockA{\textit{dept. name of organization (of Aff.)} \\
\textit{name of organization (of Aff.)}\\
City, Country \\
email address or ORCID}}

\maketitle

\begin{abstract}
This survey paper provides a comprehensive examination of the current state of multi-camera multi-object tracking, a fundamental problem in computer vision with numerous applications in surveillance, robotics, and autonomous vehicles. The paper introduces the concept of multi-object tracking, highlighting its evolution and key challenges, particularly in multi-camera settings. A thorough review of background and fundamentals of object detection and recognition is presented, emphasizing the crucial role of data association techniques, context awareness, and trajectory forecasting in achieving robust and accurate tracking systems.

The survey delves into advanced techniques, including deep learning approaches, and explores various applications of multi-object tracking, such as surveillance, autonomous vehicles, and robotics. The challenges and open problems in this field are discussed, with a focus on the need for robust algorithms that can handle differences in camera views, lighting conditions, and object appearances.

This paper contributes to the existing body of research by providing an in-depth analysis of key findings, novel methodologies, and common themes that have emerged from the literature. The scope of this survey encompasses various aspects of multi-camera multi-object tracking, including data association, context awareness, trajectory forecasting, and adversarial attacks. The contributions of this paper include a comprehensive review of the current state of multi-camera multi-object tracking, highlighting the importance of integrating multiple modalities and improving robustness in future research directions.

Overall, this survey aims to provide a thorough understanding of how multi-camera multi-object tracking works, with a focus on its key themes, findings, and contributions. The paper is intended to serve as a valuable resource for researchers and practitioners in the field of computer vision, providing insights into the current state of multi-camera multi-object tracking and highlighting future research directions.
\end{abstract}

\begin{IEEEkeywords}
large language models, multimodal learning, natural language processing, machine learning, artificial intelligence
\end{IEEEkeywords}


\section{Introduction to Multi-Object Tracking}
Introduction to Multi-Object Tracking

Multi-object tracking is a fundamental problem in computer vision, with numerous applications in surveillance, robotics, and autonomous vehicles. The task involves detecting and tracking multiple objects across frames in a video sequence, while maintaining their identities and trajectories. In recent years, the focus has shifted towards multi-camera multi-object tracking, where objects are tracked across multiple cameras with overlapping or non-overlapping fields of view \cite{The Challenge of Multi-Camera Tracking}. This introduces additional challenges, such as handling occlusions, varying lighting conditions, and inconsistent object appearances across different camera views.

According to \cite{MCTR: Multi Camera Tracking Transformer}, an end-to-end approach using transformers can effectively address these challenges by jointly detecting and tracking objects across multiple cameras. This method treats multi-camera tracking as a sequence-to-sequence problem, where the input is a sequence of frames from each camera, and the output is a sequence of object tracks. In contrast, \cite{An equalised global graphical model-based approach for multi-camera object tracking} proposes a global graph model with an improved similarity metric to handle single-camera and inter-camera tracking differently.

A common theme among recent papers is the emphasis on developing effective multi-camera tracking approaches that can handle the complexities of real-world scenarios. For instance, \cite{Tracking without bells and whistles} presents a simple yet efficient tracking approach that achieves state-of-the-art performance on three multi-object tracking benchmarks without requiring specific training or optimization on tracking data. This highlights the potential for deep learning-based methods to improve multi-object tracking performance, as also demonstrated by \cite{Tracking the Trackers: An Analysis of the State of the Art in Multiple Object Tracking}.

The use of deep learning-based approaches, such as transformers and graphical models, is a dominant pattern among recent papers. These methods have shown significant promise in improving multi-object tracking performance, particularly in scenarios with complex object interactions and occlusions \cite{MCTR: Multi Camera Tracking Transformer}. However, there are also limitations and challenges associated with these approaches, such as the need for large-scale datasets with overlapping camera views to train and evaluate their performance \cite{MCTR: Multi Camera Tracking Transformer}.

Furthermore, current benchmarks and evaluation metrics for multiple object tracking have been criticized for being inadequate, highlighting the need for more comprehensive and standardized evaluations \cite{Tracking the Trackers: An Analysis of the State of the Art in Multiple Object Tracking}. This emphasizes the importance of developing new evaluation protocols that can accurately assess the performance of multi-camera multi-object tracking methods in real-world scenarios.

In conclusion, multi-camera multi-object tracking is a complex problem that requires effective approaches to handle the challenges of tracking objects across multiple cameras with overlapping or non-overlapping fields of view. Recent papers have highlighted the potential of deep learning-based methods, such as transformers and graphical models, to improve multi-object tracking performance. However, there are also limitations and challenges associated with these approaches, emphasizing the need for further research and development in this field. As the field continues to evolve, it is essential to develop new evaluation protocols and benchmarks that can accurately assess the performance of multi-camera multi-object tracking methods, ultimately leading to more effective and efficient solutions for real-world applications \cite{The Challenge of Multi-Camera Tracking}.

\section{Background and Fundamentals of Object Detection and Recognition}
Background and Fundamentals of Object Detection and Recognition

The task of multi-camera multi-object tracking is intricately linked with object detection and recognition, as the ability to accurately detect and identify objects across multiple cameras is crucial for effective tracking \cite{Challenge of Multi-Camera Tracking}. Recent advancements in deep learning have led to significant improvements in object detection algorithms, with techniques such as end-to-end trainable models \cite{Learning to Track with Object Permanence} and transformer-based architectures \cite{MCTR: Multi Camera Tracking Transformer} demonstrating promising results. These approaches have been shown to effectively address the challenges of occlusions and maintaining object permanence in tracking, allowing for more robust and accurate detection and recognition of objects across multiple cameras.

A key theme that emerges from the literature is the importance of robust object detection and recognition algorithms \cite{Detection, Recognition, and Tracking: A Survey}. The development of novel architectures and models that can effectively reason about object locations and identities under occlusions has been a major focus of research in this area \cite{Learning to Track with Object Permanence}. For instance, the use of spatio-temporal, recurrent memory modules has been shown to be effective in maintaining object permanence and tracking objects across multiple cameras \cite{Learning to Track with Object Permanence}. Furthermore, probabilistic association of track embeddings with detections in every camera view and frame has been proposed as a means of improving the accuracy of multi-camera tracking \cite{MCTR: Multi Camera Tracking Transformer}.

The need for large-scale datasets and benchmarks to evaluate multi-camera tracking algorithms is also a common theme in the literature \cite{Challenge of Multi-Camera Tracking}. Datasets such as KITTI, MOT17, MMPTrack, and AI City Challenge have been widely used to evaluate the performance of multi-camera tracking algorithms \cite{MCTR: Multi Camera Tracking Transformer}. These datasets provide a means of comparing the performance of different algorithms and identifying areas for improvement. Moreover, simultaneous object detection, tracking, and event recognition have been explored, demonstrating the potential for multidirectional information flow between these components to improve overall performance \cite{Simultaneous Object Detection, Tracking, and Event Recognition}.

Despite the advancements in object detection and recognition, there are still several limitations and challenges associated with multi-camera tracking. The complexity of designing end-to-end trainable models that can reason about object locations and identities under occlusions is a significant challenge \cite{Learning to Track with Object Permanence}. Additionally, the requirement for large-scale datasets and benchmarks to evaluate multi-camera tracking algorithms can be a limitation, particularly in scenarios where such datasets are not readily available \cite{Challenge of Multi-Camera Tracking}. Nevertheless, the ongoing research in this area is expected to lead to further improvements in object detection and recognition, ultimately enhancing the accuracy and effectiveness of multi-camera multi-object tracking systems.

The connections between object detection, recognition, and tracking are multifaceted, with each component influencing the others. For example, accurate object detection is crucial for effective tracking, while robust tracking algorithms can provide valuable feedback to improve object detection \cite{Detection, Recognition, and Tracking: A Survey}. Furthermore, the ability to recognize objects across multiple cameras can facilitate more accurate tracking and improve overall system performance \cite{MCTR: Multi Camera Tracking Transformer}. As research in this area continues to evolve, it is likely that we will see even more sophisticated and effective multi-camera multi-object tracking systems that can accurately detect, recognize, and track objects in complex environments.

\section{Evolution of Multi-Object Tracking Algorithms}
The evolution of multi-object tracking algorithms has been marked by significant advancements in recent years, driven by the increasing demand for robust and efficient tracking systems in various applications, including surveillance, robotics, and autonomous vehicles. A key challenge in this field is the development of end-to-end methods that can handle multi-camera tracking in a unified way, addressing the complexities of inter-camera object tracking \cite{Challenge of Multi-Camera Tracking}. To this end, researchers have proposed innovative approaches, such as the Multi Camera Tracking Transformer (MCTR), which leverages DETR to produce detections and detection embeddings independently for each camera view \cite{MCTR: Multi Camera Tracking Transformer}.

The use of graph models and similarity metrics has also been explored to improve multi-camera object tracking performance. For instance, an equalised global graphical model-based approach has been proposed to optimize multi-camera object tracking by using an improved similarity metric \cite{An equalised global graphical model-based approach for multi-camera object tracking}. Furthermore, the importance of benchmarking and standardized evaluation has been emphasized, with studies highlighting the need for comprehensive evaluations to advance the field \cite{Tracking the Trackers: An Analysis of the State of the Art in Multiple Object Tracking}.

A notable trend in recent research is the development of simple yet effective tracking algorithms, such as Tracktor, which achieves state-of-the-art performance on multi-object tracking benchmarks without requiring specific training or optimization on tracking data \cite{Tracking without bells and whistles}. This approach underscores the potential for minimalist solutions to achieve robust tracking performance. In contrast, other papers have proposed more complex methods, emphasizing the importance of addressing inter-camera object tracking and the challenges posed by occlusions, missing detections, and small objects.

The analysis of recent papers reveals common themes and patterns, including the need for end-to-end methods, the importance of addressing inter-camera object tracking, and the use of deep learning techniques to improve tracking performance. The recognition of challenges such as occlusions, missing detections, and small objects is also a recurring theme \cite{Challenge of Multi-Camera Tracking}. Moreover, the emphasis on benchmarking and standardized evaluation highlights the need for comprehensive evaluations to advance the field.

Despite significant progress, multi-camera tracking remains a challenging problem due to the complexities of handling multiple cameras with overlapping fields of view. Inter-camera object tracking is particularly difficult, especially in cases where single-camera tracking is poor \cite{An equalised global graphical model-based approach for multi-camera object tracking}. The lack of standardized evaluation metrics and datasets also hinders the advancement of the field \cite{Tracking the Trackers: An Analysis of the State of the Art in Multiple Object Tracking}.

In conclusion, the evolution of multi-object tracking algorithms has been marked by significant advancements, driven by innovative approaches and a deeper understanding of the challenges involved. The development of end-to-end methods, graph models, and simple yet effective tracking algorithms has improved tracking performance, while highlighting the need for comprehensive evaluations and standardized metrics to advance the field. As research continues to address the complexities of multi-camera object tracking, we can expect to see further improvements in tracking systems, with implications for various applications that rely on robust and efficient tracking capabilities.

\section{Data Association Techniques for Tracking}
Data association techniques are a crucial component of multi-camera multi-object tracking, enabling the linking of detections across frames and cameras to form consistent trajectories. The importance of effective data association methods is highlighted in several papers, including "Challenge of Multi-Camera Tracking" \cite{Challenge of Multi-Camera Tracking}, which identifies the challenges in multi-camera tracking and suggests new directions for future work. Furthermore, "Multiple object tracking with context awareness" \cite{Multiple object tracking with context awareness} proposes the use of contextual information, such as social and spatial context, to improve data association, demonstrating the potential benefits of incorporating additional cues into the tracking process.

A key theme that emerges from the analysis of relevant papers is the need for robust data association techniques to handle challenges such as occlusions, false detections, and camera networks. Distributed approaches to multi-camera tracking, such as the one presented in "Distributed Multi-Target Tracking in Camera Networks" \cite{Distributed Multi-Target Tracking in Camera Networks}, are gaining attention, with a focus on efficient communication management and local decision making. This approach utilizes a re-identification network and a distributed tracker manager module for cross-camera data association, highlighting the potential benefits of decentralized tracking systems.

The choice of fusion method and optimization technique is also critical in data association frameworks, as demonstrated in "FusionSORT: Fusion Methods for Online Multi-object Visual Tracking" \cite{FusionSORT: Fusion Methods for Online Multi-object Visual Tracking}. This paper investigates different fusion methods for associating detections to tracklets, highlighting the importance of selecting the right approach for effective data association. Additionally, "Beyond Pixels: Leveraging Geometry and Shape Cues for Online Multi-Object Tracking" \cite{Beyond Pixels: Leveraging Geometry and Shape Cues for Online Multi-Object Tracking} introduces geometry and object shape and pose costs for multi-object tracking, demonstrating consistent improvement over state-of-the-art methods.

The use of contextual information, such as social and spatial context, is recognized as a crucial factor in improving data association. For example, "Multiple object tracking with context awareness" \cite{Multiple object tracking with context awareness} emphasizes the importance of modeling human behavior and interactions to enhance tracking performance. However, incorporating such information can be challenging due to the complexity of modeling human behavior and interactions.

Technical details and methodologies employed in data association frameworks vary across papers, including the use of Hungarian assignment, Kalman filter gating, and optimization-based methods. Features such as motion, appearance, height intersection-over-union (height-IoU), and tracklet confidence are used in different combinations to improve tracking performance. Camera networks, such as the one presented in "Distributed Multi-Target Tracking in Camera Networks" \cite{Distributed Multi-Target Tracking in Camera Networks}, require efficient communication management and local decision making, which can be complex to implement.

Despite the progress made in data association techniques for tracking, several limitations and challenges remain. Handling occlusions and false detections remains a significant challenge in multi-camera multi-object tracking. Distributed approaches require efficient communication management and local decision making, which can be difficult to achieve in practice. Choosing the right fusion method and optimization technique is crucial for effective data association, but can be challenging due to the complexity of the tracking problem.

In conclusion, data association techniques are a critical component of multi-camera multi-object tracking, and recent papers have highlighted the importance of robust methods to handle challenges such as occlusions, false detections, and camera networks. The use of contextual information, distributed approaches, and fusion methods are key themes that emerge from the analysis of relevant papers. While significant progress has been made, limitations and challenges remain, emphasizing the need for continued research in this area to develop more effective and efficient data association techniques for tracking.

\section{Context Awareness in Multiple Object Tracking}
Context awareness is a crucial aspect of multiple object tracking, particularly in scenarios involving multiple cameras. The ability to understand the context in which objects are being tracked enables more accurate and robust tracking systems. Recent research has highlighted the importance of considering context awareness in multi-object tracking, with various approaches being proposed to address the challenges associated with this task \cite{The Challenge of Multi-Camera Tracking}. 

One key finding that emerges from the analysis of relevant papers is the need for innovative techniques and system architectures that can handle occlusions, maintain object identities, and manage data association \cite{MCTR: Multi Camera Tracking Transformer}. The use of deep learning-based approaches, such as transformers and convolutional neural networks (CNNs), has been prevalent in addressing these challenges \cite{Learning to Track with Object Permanence}. For instance, the MCTR approach uses DETR to produce detections and detection embeddings independently for each camera view, which are then integrated using a set of track embeddings \cite{MCTR: Multi Camera Tracking Transformer}. 

Another significant theme that emerges from the analysis is the importance of considering hierarchical visual representations in multi-object tracking \cite{Multi-Object Tracking by Hierarchical Visual Representations}. This approach discriminates between objects using compositional visual regions, semantic cues, and background contextual information. Furthermore, the use of object permanence has been proposed as a means of allowing models to reason about object locations even under full occlusions \cite{Learning to Track with Object Permanence}. 

The analysis also reveals contrasting viewpoints and approaches, with some papers emphasizing end-to-end approaches \cite{MCTR: Multi Camera Tracking Transformer}, while others focus on traditional tracking-by-detection methods or graph-based models \cite{An equalised global graphical model-based approach for multi-camera object tracking}. The technical details and methodologies employed by each paper demonstrate the diversity of solutions being explored to address the challenges of multi-camera tracking. 

Despite the advancements in this area, several limitations and challenges remain. For instance, handling occlusions, maintaining object identities, and managing data association continue to be significant challenges \cite{The Challenge of Multi-Camera Tracking}. Additionally, the requirement for large-scale datasets for training and the potential struggles with cases where objects are not fully visible or have complex motion patterns are notable limitations \cite{MCTR: Multi Camera Tracking Transformer}. 

In conclusion, context awareness is a vital component of multiple object tracking, particularly in multi-camera scenarios. The analysis of relevant papers highlights the importance of considering innovative techniques and system architectures that can handle occlusions, maintain object identities, and manage data association. The use of deep learning-based approaches, hierarchical visual representations, and object permanence are significant themes that emerge from the analysis. While limitations and challenges remain, continued research in this area is essential to develop more effective and robust multi-object tracking systems \cite{The Challenge of Multi-Camera Tracking}. Ultimately, the development of context-aware multi-object tracking systems has significant implications for various applications, including surveillance, robotics, and autonomous vehicles, where accurate and robust tracking is critical.

\section{Trajectory Forecasting and Long-Term Tracking}
Trajectory forecasting and long-term tracking are critical components of multi-camera multi-object tracking systems, enabling the prediction of future locations of objects across different camera views and the maintenance of consistent tracks over time. Recent research has made significant strides in these areas, with several key findings and contributions standing out. For instance, the concept of trajectory tensors for multi-camera trajectory forecasting has been introduced, allowing for the prediction of future locations in all possible viewpoints and outperforming existing single-camera methods \cite{Multi-Camera Trajectory Forecasting with Trajectory Tensors}. Additionally, end-to-end frameworks such as "Past-and-Future reasoning for Tracking" (PF-Track) have been proposed, integrating past and future reasoning to significantly improve tracking performance on datasets like nuScenes \cite{Standing Between Past and Future: Spatio-Temporal Modeling for Multi-Camera 3D Multi-Object Tracking}.

The importance of developing methods that can handle multi-camera setups is a common theme throughout the literature, with single-camera approaches being limited in their ability to track objects across non-overlapping views \cite{Challenge of Multi-Camera Tracking}. Spatio-temporal modeling has also emerged as a crucial aspect of trajectory forecasting and long-term tracking, with techniques such as transformers and attention mechanisms being leveraged to incorporate both spatial and temporal information \cite{MCTR: Multi Camera Tracking Transformer}. The use of advanced mathematical and computational techniques is another notable trend, reflecting the complexity of the problem and the need for sophisticated models to achieve accurate results.

The shift from single-camera to multi-camera tracking has introduced unique challenges, such as camera synchronization and object re-identification across views. Researchers have proposed various solutions to address these challenges, including the use of trajectory tensors \cite{Multi-Camera Trajectory Forecasting with Trajectory Tensors} and probabilistic associations \cite{MCTR: Multi Camera Tracking Transformer}. The development of new datasets, such as the Warwick-NTU Multi-camera Forecasting Database \cite{Multi-Camera Trajectory Forecasting: Pedestrian Trajectory Prediction in a Network of Cameras}, has also facilitated research in this area by providing valuable resources for evaluating and improving multi-camera tracking algorithms.

Despite these advancements, several limitations and challenges remain. Scalability and efficiency are significant concerns, as complex models may face challenges in terms of computational efficiency and scalability to larger camera networks or longer tracking durations \cite{MCTR: Multi Camera Tracking Transformer}. Dealing with long-term occlusions and accurately re-identifying objects across non-overlapping views also remain technical hurdles \cite{Challenge of Multi-Camera Tracking}. Furthermore, there is a need for more diverse and large-scale datasets to comprehensively evaluate and improve multi-camera tracking algorithms \cite{Multi-Camera Trajectory Forecasting: Pedestrian Trajectory Prediction in a Network of Cameras}.

In conclusion, trajectory forecasting and long-term tracking are essential components of multi-camera multi-object tracking systems, and recent research has made significant progress in these areas. The development of novel techniques, such as trajectory tensors and end-to-end frameworks, has improved the accuracy and efficiency of multi-camera tracking algorithms. However, challenges remain, and further research is needed to address issues like scalability, occlusions, and re-identification. As the field continues to evolve, it is likely that we will see the development of even more sophisticated models and techniques, enabling the creation of robust and efficient multi-camera tracking systems.

\section{Adversarial Attacks on Multi-Camera Systems for Object Tracking}
The vulnerability of multi-camera tracking systems to adversarial attacks has become a growing concern in recent years, with significant implications for various applications, including surveillance and autonomous vehicles \cite{MCTR}. Developing robust systems that can withstand such attacks is crucial, as highlighted by several studies that have investigated the challenges and vulnerabilities of multi-camera tracking systems \cite{Physical Adversarial Examples for Multi-Camera Systems}. For instance, MCTR introduces a novel end-to-end approach for multi-object detection and tracking across multiple cameras with overlapping fields of view, demonstrating the potential for improved robustness against adversarial attacks \cite{MCTR}. In contrast, Physical Adversarial Examples for Multi-Camera Systems proposes a new attack method called Transcender-MC, which incorporates online 3D renderings and perspective projections to optimize attacks on multiple perspectives at once, highlighting the ongoing cat-and-mouse game between attackers and defenders in the field of object tracking \cite{Physical Adversarial Examples for Multi-Camera Systems}.

The analysis of relevant papers reveals several key themes and developments in the area of adversarial attacks on multi-camera systems for object tracking. Firstly, the importance of developing robust multi-camera tracking systems that can withstand adversarial attacks is a recurring theme, with studies emphasizing the need for end-to-end approaches that can handle multiple cameras with overlapping fields of view \cite{MCTR}. Secondly, the vulnerability of neural networks to physical adversarial examples, particularly in air-gapped scenarios, is a significant concern, as highlighted by Physical Adversarial Examples for Multi-Camera Systems \cite{Physical Adversarial Examples for Multi-Camera Systems}. Thirdly, the significance of evaluating the effectiveness of existing attack methods on various object trackers and datasets is underscored by studies such as Reproducibility Study on Adversarial Attacks Against Robust Transformer Trackers, which evaluates the effectiveness of existing adversarial attacks on object trackers with transformer and non-transformer backbones \cite{Reproducibility Study on Adversarial Attacks Against Robust Transformer Trackers}.

The technical details and methodologies employed in these studies are also noteworthy. For example, MCTR uses DEtector TRansformer (DETR) to produce detections and detection embeddings independently for each camera view, while Physical Adversarial Examples for Multi-Camera Systems incorporates online 3D renderings and perspective projections in the training process \cite{MCTR}. An equalised global graphical model-based approach, on the other hand, uses an improved similarity metric to treat similarities between single camera tracking and inter-camera tracking differently \cite{An equalised global graphical model-based approach}. These technical details highlight the complexity of multi-camera systems, including overlapping fields of view and varying camera perspectives, which poses significant technical challenges for developing robust tracking systems.

Despite these challenges, the development of robust multi-camera tracking systems that can withstand adversarial attacks is crucial for various applications. The implications of adversarial attacks on multi-camera systems are far-reaching, with potential consequences for surveillance, autonomous vehicles, and other safety-critical applications \cite{Reproducibility Study on Adversarial Attacks Against Robust Transformer Trackers}. Therefore, continued research in this area is essential to address the limitations and challenges highlighted by existing studies. The connections between adversarial attacks, object tracking, and computer vision are also noteworthy, as they underscore the need for a comprehensive approach that considers the interplay between these areas \cite{Physical Adversarial Examples for Multi-Camera Systems}. Ultimately, the development of robust multi-camera tracking systems that can withstand adversarial attacks requires a deep understanding of the challenges and vulnerabilities of these systems, as well as the ongoing efforts to address them.

\section{Advanced Techniques: Deep Learning for Multi-Object Tracking}
Advanced Techniques: Deep Learning for Multi-Object Tracking

The advent of deep learning has significantly impacted the field of multi-object tracking (MOT), enabling the development of more accurate and efficient tracking systems. As highlighted in "Deep Learning in Video Multi-Object Tracking: A Survey" \cite{Deep Learning in Video Multi-Object Tracking: A Survey}, deep learning models have been employed in various stages of MOT algorithms, leveraging their representational power to improve object detection, feature extraction, and association. For instance, the use of convolutional neural networks (CNNs) for object detection has become a standard approach, as seen in \cite{End-to-End Multi-Object Tracking with Global Response Map}. Moreover, the integration of deep learning with traditional methods, such as optical flow algorithms, has been explored in "A Hybrid Approach To Real-Time Multi-Object Tracking" \cite{A Hybrid Approach To Real-Time Multi-Object Tracking}, achieving a balance between performance and computational costs.

The application of deep learning to MOT has also led to the development of end-to-end approaches, which simplify the tracking process by integrating detection, association, and tracking into a unified framework. For example, "End-to-End Multi-Object Tracking with Global Response Map" \cite{End-to-End Multi-Object Tracking with Global Response Map} presents an end-to-end approach using a global response map generated over frames, allowing for accurate object tracking without the need for explicit object detection and association steps. Similarly, "MCTR: Multi Camera Tracking Transformer" \cite{MCTR: Multi Camera Tracking Transformer} proposes an end-to-end multi-camera tracking framework using transformers, enabling consistent object tracks across multiple cameras with overlapping fields of view through probabilistic associations.

A common theme among these papers is the challenge of balancing computational costs and tracking performance, particularly in real-time applications. The use of deep learning models can significantly improve accuracy but often at the cost of increased computational requirements. To address this issue, researchers have explored hybrid approaches, such as combining traditional methods with deep learning, or optimizing deep learning architectures for efficient processing. For instance, \cite{A Hybrid Approach To Real-Time Multi-Object Tracking} demonstrates a hybrid strategy that combines classical optical flow algorithms with deep learning architectures for real-time tracking.

The expansion of MOT applications to multi-camera scenarios has also been recognized as a critical area for advancement. Papers like "MCTR: Multi Camera Tracking Transformer" \cite{MCTR: Multi Camera Tracking Transformer} highlight the potential of comprehensive surveillance and monitoring using multiple cameras. However, this also introduces new challenges, such as handling occlusions, appearance changes, and diverse scenarios in video surveillance. To overcome these limitations, researchers have proposed novel deep learning architectures, such as transformers, which can effectively handle multi-camera inputs and produce consistent tracks across views.

Despite the significant progress made in applying deep learning to MOT, several challenges persist. One of the primary concerns is the computational cost associated with deep learning models, which can limit their applicability in real-time applications. Additionally, handling object occlusions, appearance changes, and diverse scenarios in video surveillance remains a significant challenge. The development of large-scale datasets for training and evaluating deep learning models in MOT is also essential, as highlighted by experiments on datasets like MMPTrack and AI City Challenge.

In conclusion, the application of deep learning to multi-object tracking has led to significant advancements in the field, with the development of more accurate and efficient tracking systems. While challenges remain, the progress made in addressing these issues indicates a promising future for MOT systems in various applications. The importance of continued research into end-to-end approaches, hybrid models, and novel deep learning architectures, such as transformers, cannot be overstated, as these developments have the potential to overcome existing limitations and improve tracking performance. As highlighted in "Deep Learning in Video Multi-Object Tracking: A Survey" \cite{Deep Learning in Video Multi-Object Tracking: A Survey}, the future of MOT lies in the effective integration of deep learning with traditional methods, enabling the development of more robust and efficient tracking systems.

\section{Applications of Multi-Object Tracking: Surveillance, Autonomous Vehicles, and Beyond}
The applications of multi-object tracking have expanded significantly in recent years, with surveillance, autonomous vehicles, and robotics being some of the most prominent areas of research. As highlighted in \cite{Challenge of Multi-Camera Tracking}, the challenge of maintaining consistent object tracks across multiple cameras and frames is a common problem that plagues many existing algorithms. However, researchers have made significant strides in addressing this issue, with papers like \cite{No Blind Spots: Full-Surround Multi-Object Tracking for Autonomous Vehicles} presenting modular frameworks for tracking multiple objects using cameras and LiDARs. This approach enables the production of continuous object tracks across multiple sensors and modalities, which is crucial for applications such as autonomous vehicles.

The importance of efficient and accurate multi-camera tracking in various applications is a recurring theme across many papers, including \cite{Tracking the Trackers: An Analysis of the State of the Art in Multiple Object Tracking} and \cite{CORT: Class-Oriented Real-time Tracking for Embedded Systems}. The use of sensor fusion, such as combining camera and LiDAR data, is also a common pattern in several papers, including \cite{No Blind Spots} and \cite{MCTR: Multi Camera Tracking Transformer}. This approach enables the leveraging of complementary strengths of different sensors to improve tracking performance. Furthermore, the need for real-time performance and low latency is highlighted in multiple papers, including \cite{CORT} and \cite{MCTR}, as this is critical for many applications, particularly those that involve autonomous systems.

The technical details of multi-object tracking algorithms vary widely across different papers, with some proposing class-oriented approaches like \cite{CORT}, while others introduce end-to-end transformer-based approaches like \cite{MCTR}. The former splits the Hungarian matrix by class and invokes the second re-identification stage only when necessary, reducing execution times without penalizing tracking performance. In contrast, the latter leverages end-to-end detectors like DEtector TRansformer (DETR) to produce detections and detection embeddings independently for each camera view. Additionally, papers like \cite{Tracking the Trackers} provide benchmarking analyses of state-of-the-art trackers, which helps to identify current trends and weaknesses in multiple people tracking methods.

Despite the significant progress made in multi-object tracking research, several limitations and challenges remain, including maintaining consistent object tracks across multiple cameras and frames, handling occlusions and blind spots, achieving real-time performance and low latency, integrating data from different sensor modalities, and addressing the complexity of urban scenarios with multiple objects and classes. As highlighted in \cite{Challenge of Multi-Camera Tracking}, these challenges are particularly pronounced in multi-camera tracking applications, where the presence of multiple cameras and sensors can introduce additional complexities.

In conclusion, the applications of multi-object tracking have expanded significantly in recent years, with surveillance, autonomous vehicles, and robotics being some of the most prominent areas of research. While researchers have made significant strides in addressing the challenges associated with multi-object tracking, several limitations and challenges remain. The development of more efficient, accurate, and robust multi-object tracking algorithms that can handle the complexities of real-world scenarios is crucial for the advancement of many applications, particularly those that involve autonomous systems. As the field continues to evolve, it is likely that we will see significant advancements in areas like sensor fusion, real-time performance, and class-oriented approaches, which will help to address the current limitations and challenges in multi-object tracking research.

\section{Challenges and Open Problems in Multi-Object Tracking}
The field of multi-object tracking has witnessed significant advancements in recent years, driven by the increasing demand for robust and efficient tracking algorithms in various applications. However, despite these advances, several challenges and open problems remain to be addressed. As highlighted in \cite{Challenge of Multi-Camera Tracking}, multi-camera tracking poses significant technological and system architecture challenges, including the need for robust algorithms that can handle different camera views and overlapping fields of view.

One of the key findings in this area is the importance of global graphical models in improving multi-camera object tracking performance. For instance, \cite{An equalised global graphical model-based approach for multi-camera object tracking} introduces a global graph model with an improved similarity metric, demonstrating its effectiveness in enhancing overall tracking performance. In contrast, end-to-end approaches have also gained popularity, with papers like \cite{MCTR: Multi Camera Tracking Transformer} proposing novel architectures that leverage detectors like DETR to produce detections and detection embeddings independently for each camera view.

The evaluation and benchmarking of multi-object tracking algorithms are also crucial aspects that require attention. As noted in \cite{Tracking the Trackers: An Analysis of the State of the Art in Multiple Object Tracking}, current evaluation metrics may not accurately reflect real-world performance, highlighting the need for more comprehensive benchmarks. Furthermore, the trend towards end-to-end approaches is evident, with papers like \cite{MCTR: Multi Camera Tracking Transformer} and \cite{Tracking without bells and whistles} demonstrating the potential of these methods in simplifying tracking pipelines and improving performance.

A common theme that emerges from the analyzed papers is the need for robust multi-camera tracking algorithms that can handle different camera views and overlapping fields of view. This is echoed in \cite{Challenge of Multi-Camera Tracking}, which emphasizes the importance of developing algorithms that can address these challenges. Additionally, the importance of evaluation and benchmarking is highlighted in \cite{Tracking the Trackers: An Analysis of the State of the Art in Multiple Object Tracking}, which presents a benchmark for multiple object tracking and provides an in-depth analysis of state-of-the-art trackers.

The technical details and methodologies employed in the analyzed papers are diverse, ranging from graphical models to transformers and object detectors. For example, \cite{An equalised global graphical model-based approach for multi-camera object tracking} uses a global graph model with an improved similarity metric, while \cite{MCTR: Multi Camera Tracking Transformer} leverages transformers like DETR to produce detections and detection embeddings independently for each camera view. In contrast, \cite{Tracking without bells and whistles} exploits the bounding box regression of an object detector to predict the position of an object in the next frame.

Despite the advancements in multi-object tracking, several limitations and challenges remain. Scalability and robustness are significant concerns, particularly in scenes with multiple objects, occlusions, or varying lighting conditions. As noted in \cite{Challenge of Multi-Camera Tracking}, current algorithms often struggle with these challenges, highlighting the need for more robust and efficient methods. Furthermore, the development of algorithms that can handle real-world applications, such as varying camera views and overlapping fields of view, is crucial. This is emphasized in papers like \cite{MCTR: Multi Camera Tracking Transformer} and \cite{Challenge of Multi-Camera Tracking}, which highlight the importance of developing practical and effective solutions for real-world scenarios.

In conclusion, the challenges and open problems in multi-object tracking are multifaceted and require a comprehensive approach to address. The development of robust and efficient algorithms that can handle different camera views, overlapping fields of view, and varying lighting conditions is crucial. Furthermore, the evaluation and benchmarking of these algorithms must be rigorous and comprehensive, taking into account real-world performance and applications. By addressing these challenges and limitations, researchers can develop more effective and practical solutions for multi-object tracking, ultimately driving advancements in various fields, including computer vision, robotics, and surveillance.

\section{Future Directions: Integrating Multiple Modalities and Improving Robustness}
Future directions in multi-camera multi-object tracking are poised to focus on integrating multiple modalities and improving robustness, driven by the need for more accurate and reliable tracking systems in various applications. The analysis of relevant papers reveals a trend towards leveraging different sensor data, such as camera and LiDAR, to enhance tracking performance \cite{CrossTracker}. This multi-modal approach is exemplified by CrossTracker's two-stage paradigm, which utilizes a multi-modal modeling module for robust metric calculation and a trajectory refinement module for cross correction between camera and LiDAR streams \cite{CrossTracker}. Similarly, MCTR proposes an end-to-end approach using transformers to integrate detection embeddings and track embeddings for probabilistic association \cite{MCTR}.

The integration of multiple modalities is expected to play a crucial role in improving tracking robustness, particularly in scenarios with occlusions or missing data. For instance, PF-Track's spatio-temporal modeling with past and future reasoning demonstrates the effectiveness of incorporating temporal information and spatial relationships between objects \cite{PF-Track}. This approach reduces ID-Switches by 90\% on the nuScenes dataset, highlighting the potential benefits of such methods. Furthermore, graph-based models, such as the equalised global graphical model-based approach, offer a promising direction for representing and reasoning about object tracks across multiple cameras \cite{GlobalGraphModel}.

The shift towards end-to-end approaches is another notable development in multi-camera tracking. Methods like MCTR and PF-Track represent a departure from traditional heuristic techniques, instead opting for differentiable losses and end-to-end training \cite{MCTR, PF-Track}. This trend is expected to continue, with future research focusing on developing more sophisticated end-to-end architectures that can learn to optimize tracking performance directly from data. The use of transformers, as seen in MCTR, is particularly noteworthy, as it enables the integration of detection embeddings and track embeddings for probabilistic association \cite{MCTR}.

Despite these advancements, several challenges remain, including scalability, occlusions, and missing data. Multi-camera tracking systems can be computationally expensive and require significant resources, making scalability a pressing concern \cite{Scalability}. Additionally, handling long-term occlusions and missing data remains a challenge, with most methods relying on heuristic techniques or simplistic assumptions \cite{Occlusions}. The effective combination of different sensor data is also an ongoing research challenge, with many approaches struggling to balance the contributions of each modality \cite{ModalitiesIntegration}.

To address these challenges, future research should focus on developing more efficient and robust tracking algorithms that can handle complex scenarios. This may involve exploring novel architectures, such as those leveraging attention mechanisms or graph neural networks, to better model spatial and temporal relationships between objects \cite{AttentionMechanisms, GraphNeuralNetworks}. Moreover, the development of robust evaluation metrics for multi-camera tracking performance is essential to compare different approaches and drive progress in the field \cite{EvaluationMetrics}.

In conclusion, the future of multi-camera multi-object tracking lies in integrating multiple modalities and improving robustness. By leveraging different sensor data, developing end-to-end architectures, and exploring novel modeling techniques, researchers can create more accurate and reliable tracking systems. Addressing the challenges of scalability, occlusions, and missing data will be crucial to realizing the full potential of these systems, with implications for various applications, including surveillance, autonomous vehicles, and smart cities \cite{Surveillance, AutonomousVehicles, SmartCities}. As the field continues to evolve, it is likely that we will see significant advancements in multi-camera tracking, driven by the increasing availability of large-scale datasets and advances in computing hardware \cite{LargeScaleDatasets, ComputingHardware}.

\section{Conclusion and Summary of Key Findings in Multi-Object Tracking Research}
In conclusion, this survey of multi-object tracking research highlights several key findings and contributions that have significantly advanced our understanding of how multi-camera multi-object tracking works. The challenge of multi-camera tracking has been emphasized in various studies, including "The Challenge of Multi-Camera Tracking" \cite{The Challenge of Multi-Camera Tracking}, which underscores the differences between single-camera and multi-camera tracking, introducing new technological and system architecture challenges. Furthermore, benchmarks for evaluating multiple object tracking methods have been proposed, such as in "Tracking the Trackers" \cite{Tracking the Trackers}, providing an analysis of state-of-the-art trackers and identifying current trends and weaknesses in people tracking methods.

The development of end-to-end approaches, like the one presented in "MCTR: Multi Camera Tracking Transformer" \cite{MCTR: Multi Camera Tracking Transformer}, leverages transformers to produce detections and track embeddings, offering a promising direction for multi-object detection and tracking across multiple cameras. Additionally, global graph models with improved similarity metrics have been proposed, as seen in "An equalised global graphical model-based approach for multi-camera object tracking" \cite{An equalised global graphical model-based approach for multi-camera object tracking}, treating single-camera and inter-camera tracking similarities differently.

Common themes and patterns emerge from the analysis of these papers, including the acknowledgement of unique challenges posed by multi-camera tracking, such as handling non-overlapping fields of view, varying camera settings, and coordinating tracks across cameras. The need for standardized evaluation metrics and benchmarks is also a recurring theme, with "Multiple Object Tracking: A Literature Review" \cite{Multiple Object Tracking: A Literature Review} providing a comprehensive review of the MOT problem and emphasizing the importance of future research directions.

Contrasting viewpoints or approaches are also evident, such as the distinction between heuristic techniques and end-to-end methods, with the latter offering potential advantages in terms of performance and efficiency. The treatment of single-camera and inter-camera tracking as distinct problems is another area of contrast, highlighting the need for tailored solutions that address the specific challenges of each scenario.

From a technical standpoint, transformer architectures and graphical models have been employed to tackle multi-object tracking challenges, with "MCTR: Multi Camera Tracking Transformer" \cite{MCTR: Multi Camera Tracking Transformer} leveraging transformers like DETR to produce detections and track embeddings. Benchmarking efforts, such as those presented in "Tracking the Trackers" \cite{Tracking the Trackers}, have tested state-of-the-art trackers on extensive datasets, providing valuable insights into their strengths and weaknesses.

Despite these advances, limitations and challenges persist, including scalability and complexity issues associated with multi-camera tracking systems, occlusions and appearance changes that can significantly impact tracking performance, and the ongoing need for standardized datasets, evaluation metrics, and benchmarks. Addressing these challenges will be crucial for advancing the state-of-the-art in multi-object tracking and enabling more effective and efficient solutions for real-world applications.

In summary, this survey highlights the key findings, contributions, and developments in multi-object tracking research, emphasizing the importance of standardized evaluation, end-to-end approaches, and tailored solutions for distinct tracking scenarios. As the field continues to evolve, it is essential to address the limitations and challenges that remain, driving progress toward more accurate, efficient, and scalable multi-object tracking systems. The connections between these developments underscore the complex and multifaceted nature of multi-camera multi-object tracking, requiring a comprehensive understanding of the underlying principles, challenges, and opportunities. Ultimately, this research has significant implications for various applications, including surveillance, autonomous vehicles, and smart cities, where accurate and efficient tracking of multiple objects is critical for safety, security, and decision-making.

\section{Conclusion}
In conclusion, this survey provides an in-depth examination of the current state of multi-camera multi-object tracking, highlighting key findings, novel methodologies, and common themes that have emerged from the analysis of relevant papers \cite{An equalised global graphical model-based approach}, \cite{MCTR: Multi Camera Tracking Transformer}, \cite{Tracking without bells and whistles}, and \cite{Tracktor}. The collective research demonstrates improved performance in multi-camera object tracking through various approaches, including global graph models, end-to-end methods, and simple yet effective trackers. Novel methodologies, such as the use of transformers for multi-camera tracking \cite{MCTR: Multi Camera Tracking Transformer} and the exploitation of object detectors for tracking \cite{Tracktor}, have been proposed to address the unique challenges posed by multi-camera tracking.

The importance of considering global information when performing multi-camera tracking is a common theme among the papers, with several highlighting the need to handle non-overlapping fields of view and inter-camera object tracking \cite{An equalised global graphical model-based approach}. End-to-end approaches have also gained popularity, as seen in \cite{MCTR: Multi Camera Tracking Transformer} and \cite{Tracktor}, which provide better performance compared to heuristic techniques. However, the complexity of tracking methods is a topic of debate, with some papers arguing that simple approaches can be effective \cite{Tracking without bells and whistles}, while others propose more complex methodologies.

The analysis also reveals contrasting viewpoints on the role of training data in multi-camera tracking, with some papers emphasizing its importance \cite{MCTR: Multi Camera Tracking Transformer} and others demonstrating good performance without specific training or optimization on tracking data \cite{Tracktor}. Technical details and methodologies, such as graph models \cite{An equalised global graphical model-based approach}, transformers \cite{MCTR: Multi Camera Tracking Transformer}, and object detectors \cite{Tracktor}, have been discussed in the context of multi-camera tracking.

Despite significant progress, limitations and challenges persist, including scalability and efficiency issues, occlusions and missing detections, and the need for standardized benchmarks \cite{Tracking the Trackers}. Future research directions should focus on developing efficient, scalable, and robust methods that can handle complex scenarios and provide improved performance. The development of standardized benchmarks is crucial for evaluating multi-camera tracking methods, but there may be limitations and biases in current benchmarking approaches. By addressing these challenges and building upon existing research, the field of multi-camera multi-object tracking can continue to advance and provide effective solutions for real-world applications. Ultimately, this survey highlights the key themes and developments in multi-camera multi-object tracking, discussing implications and connections between different approaches and methodologies, and providing a foundation for future research in this area \cite{MCTR: Multi Camera Tracking Transformer}, \cite{Tracktor}.

\end{document}